\section*{第 \ref{ch:b1:first-law} 章——热力学第一定律}

\begin{solution}{exo:b1:first-law:1}
$Q = C_{P,m}\Delta T = 29.1 \times 100 = \qty{2.91}{kJ}$; $W = -nR\Delta T =
\qty{-0.83}{kJ}$; $\Delta U = Q + W = \qty{2.08}{kJ} = C_{V,m}\Delta T$.
\end{solution}

\begin{solution}{exo:b1:first-law:2}
$W = -nRT\ln(V_2/V_1) = -2494\ln(0.25) = \qty{+3.46}{kJ}$；由 $\Delta U = 0$ 得
$Q = \qty{-3.46}{kJ}$（传给\omterm{def:b1:first-law:heat}{热源}）。
\end{solution}

\begin{solution}{exo:b1:first-law:3}
$T_2 = T_1 \times 10^{\gamma - 1} = 300 \times 10^{0.4} = \qty{754}{K}$; $P_2 =
10^{1.4} = \qty{25}{bar}$.
\end{solution}

\begin{solution}{exo:b1:first-law:4}
氩：$C_{V,m} = \tfrac32 R = 12.5$，$C_{P,m} = \tfrac52 R = 20.8$ J/(mol K)，
$\gamma = 1.67$；氮气：$20.8$、$29.1$、$\gamma = 1.40$。对于水，
加热时膨胀很小，因此定压\omterm{def:b1:first-law:transformation}{过程}中的功 $P\,\dd V$
相对于热量可忽略：$c_P - c_V \approx 0$。
\end{solution}

\begin{solution}{exo:b1:first-law:5}
$T_f = (200 \times 80 + 300 \times 20)/500 = \qty{44}{\degreeCelsius}$。计入杯子后：
$(200 \times 4.18 \times 80 + (300 \times 4.18 + 80) \times 20)/(500 \times
4.18 + 80) = \qty{43}{\degreeCelsius}$。
\end{solution}

\begin{solution}{exo:b1:first-law:6}
冰升温至 $0$：$0.050 \times 2100 \times 10 = \qty{1.05}{kJ}$；熔化：
\qty{16.7}{kJ}；合计 \qty{17.8}{kJ}。水降至
$0$ 时可放出 \qty{25.1}{kJ}，冰全部熔化；随后 $0.250 \times 4180\,T_f = 25.1 - 17.8$ kJ，得 $T_f =
\qty{7}{\degreeCelsius}$。若取 \qty{150}{g} 冰，则需 \qty{53.3}{kJ} $>$ \qty{25.1}{kJ}：
只有 $(25.1 - 3.15)/334 = \qty{66}{g}$ 的冰熔化，$T_f = \qty{0}{\degreeCelsius}$，
仍余 \qty{84}{g} 冰。
\end{solution}

\begin{solution}{exo:b1:first-law:7}
$W = -P_1(V_2 - V_1) + 0 - P_2(V_1 - V_2) + 0 = (P_2 - P_1)(V_2 - V_1) > 0$:
循环沿逆时针方向进行（低压下膨胀，
高压下压缩），因此气体\emph{得到}净功，其数值等于
所围面积；若沿顺时针方向进行，则对外输出净功——这就是热机。
\end{solution}

\begin{solution}{exo:b1:first-law:8}
$T_2 = 300 \times 7^{0.286} = \qty{523}{K}$。压缩后变热的空气
每次冲程都通过传导加热气筒。轮胎中的空气在定容下冷却至
环境温度，此时 $P \propto T$，压强按温度之比
下降——需重新充气。
\end{solution}

\begin{solution}{exo:b1:first-law:9}
$W = 0$（无外压可抗衡），$Q = 0$：$\Delta U = 0$，
因此对\omterm{def:b1:kinetic-theory:model}{理想气体} $\Delta T = 0$。\omterm{rem:b1:first-law:irreversible}{不可逆过程}：气体不会
自行流回原处。实际气体的分子间相互吸引，会略微降温
（随着分子彼此分开，部分\omterm{thm:b1:work-and-energy:ke}{动能}转化为\omterm{def:b1:work-and-energy:conservative}{势能}）。
\end{solution}

\begin{solution}{exo:b1:first-law:10}
$Q = L_v = \qty{2.26}{MJ}$; $W = -P(V_g - V_l) = -1.013\times10^5 \times 1.67 =
\qty{-0.17}{MJ}$；$\Delta U = \qty{2.09}{MJ}$：热量的 $93\%$ 用于
使分子分开（\omterm{def:b1:kinetic-theory:U}{内能}），$7\%$ 用于
克服大气压做功。
\end{solution}

\begin{solution}{exo:b1:first-law:11}
$W = -P_2(V_2 - V_1)$, $\Delta U = nC_{V,m}(T_2 - T_1)$, $V_2 = nRT_2/P_2$, $V_1
= nRT_1/P_1$: $\tfrac52(T_2 - 300) = -T_2 + 300 \times 5$, $3.5T_2 = 2250$,
$T_2 = \qty{643}{K}$，$V_2 = \qty{10.7}{L}$。可逆时：$T_2 = 300 \times
5^{0.286} = \qty{475}{K}$，$V_2 = \qty{7.9}{L}$。突然压缩对气体所做的
功更多（整个 \qty{5}{bar} 的压强从一开始就作用），末态
温度更高，压缩得更少。
\end{solution}

\begin{solution}{exo:b1:first-law:12}
$\rho = PM/RT = \qty{1.20}{kg/m^3}$: $\sqrt{\gamma P/\rho} = \sqrt{1.4 \times
1.013\times10^5/1.20} = \qty{343}{m/s}$; $\sqrt{P/\rho} = \qty{290}{m/s}$, $15\%$
偏低。由 $P/\rho = RT/M$ 得 $c = \sqrt{\gamma RT/M}$；与 $u = \sqrt{3RT/M}$ 相比
= \qty{502}{m/s}: $c/u = \sqrt{\gamma/3} = 0.68$.
\end{solution}

\begin{solution}{pb:b1:first-law:1}
\textbf{1.} $n = P_1V_1/RT_1 = 10^5 \times 2.12\times10^{-4}/2494 = \qty{8.5e-3}{mol}$.

\textbf{2.} $PV^\gamma$ 为常量：$V_2 = V_1(P_1/P_2)^{1/\gamma} = 212 \times
7^{-0.714} = \qty{53}{cm^3}$.

\textbf{3.} $T_2 = T_1(P_2/P_1)^{(\gamma - 1)/\gamma} = 300 \times 7^{0.286} =
\qty{523}{K}$.

\textbf{4.} $W = nC_{V,m}(T_2 - T_1) = 8.5\times10^{-3} \times 20.8 \times 223 =
\qty{39}{J}$.

\textbf{5.} 需补充的空气：$\Delta n = \Delta P\,V/RT = 6\times10^5 \times 2\times10^{-3}/
2494 = \qty{0.48}{mol}$：$0.48/8.5\times10^{-3} = 57$ 次冲程，约
\qty{2.2}{kJ}——相当于一分钟实实在在的体力劳动。

\textbf{6.} 定容时 $P \propto T$：空气冷却，压强随之下降
（对最后一次冲程的空气，至多降为 $300/523$ 倍）；骑车人
停顿后补打几次。

\textbf{7.} $W_{\text{等温}} = nRT\ln(P_2/P_1) = 8.5\times10^{-3} \times 2494 \times
1.95 = \qty{41}{J} > \qty{39}{J}$：\omterm{prop:b1:first-law:transformations}{等温过程}使气体
进一步压缩（体积为 \qty{30}{cm^3}，而非 \qty{53}{cm^3}）以达到 \qty{7}{bar}，
因为温度保持较低——扫过的体积更大，做功也就更多。

\textbf{8.} $n = 10^5 \times 5\times10^{-4}/2494 = \qty{0.020}{mol}$; $T_2 =
300 \times 20^{0.4} = \qty{994}{K}$.

\textbf{9.} $P_2 = 20^{1.4} = \qty{66}{bar}$.

\textbf{10.} \qty{994}{K} 远高于燃料的自燃点：喷入的燃油
接触即着火。汽油机只将其空气–燃油混合气压缩到约
\qty{750}{K}，以免在火花塞点火之前自行着火；
压缩比过高则会产生“爆震”。

\textbf{11.} $W = nC_{V,m}(T_2 - T_1) = 0.020 \times 20.8 \times 694 = \qty{289}{J}$.

\textbf{12.} $25 \times 289 = \qty{7.2}{kW}$，此为每缸的数值，
能量储存于高温气体中，并在膨胀\omterm{def:b1:first-law:transformation}{过程}中返还。

\textbf{13.} 不可逆性（气体不沿\omterm{def:b1:first-law:transformation}{准静态
过程}进行，多余的功被耗散）使 $T_2$ 升高；传给器壁的热量
使其降低。在真实热机中通常后者占优：$T_2$ 略低于
理想值。

\textbf{14.} $Q = 20\times10^{-6} \times 43\times10^6 = \qty{860}{J}$.

\textbf{15.} 等压\omterm{def:b1:first-law:transformation}{过程}：$Q = nC_{P,m}(T_3 - T_2)$：$T_3 - T_2 = 860/(0.020 \times
29.1) = \qty{1480}{K}$，$T_3 = \qty{2470}{K}$；$V_3 = V_2T_3/T_2 = 25 \times
2.48 = \qty{62}{cm^3}$.

\textbf{16.} $W = -P_2(V_3 - V_2) = -66\times10^5 \times 37\times10^{-6} =
\qty{-244}{J}$：气体对外做的功。

\textbf{17.} $T_4 = T_3(V_3/V_1)^{\gamma - 1} = 2470 \times (0.124)^{0.4} =
\qty{1070}{K}$; $W = nC_{V,m}(T_4 - T_3) = 0.020 \times 20.8 \times (-1400) =
\qty{-582}{J}$.

\textbf{18.} 净输出功 $= 244 + 582 - 289 = \qty{537}{J}$；比值
$537/860 = 62\%$（这些参数下的理想柴油循环效率；真实
热机约为 $40\%$）。

\textbf{19.} $1 - 300/2470 = 88\%$：理想柴油循环仍
远低于卡诺上限，因为它的吸热是在一段温度范围内完成的，
而非仅在最高温度下完成。

\textbf{20.} 一个循环中 $\Delta U = 0$：$Q_{\text{放出}} = -(860 - 537) =
\qty{-323}{J}$，即\ \qty{323}{J} 随废气排出；校验：
$nC_{V,m}(T_4 - T_1) = 0.020 \times 20.8 \times 770 = \qty{320}{J}$.

\textbf{21.} 振荡很快（仅需一秒），相比通过 \qty{10}{L} 空气的热传导，
故为\omterm{prop:b1:first-law:transformations}{绝热过程}。线性化
$PV^\gamma$：$\dd P/P = -\gamma\,\dd V/V = -\gamma Sx/V$。

\textbf{22.} 作用于小球的力为 $S\,\dd P = -\gamma PS^2x/V$：$m\ddot x =
-(\gamma PS^2/V)x$，为\omterm{prop:b1:point-kinematics:harmonic}{简谐运动}，$\omega_0^2 = \gamma PS^2/mV$。

\textbf{23.} $\omega_0^2 = 1.4 \times 10^5 \times 4\times10^{-8}/(0.0165 \times 0.010)
= \qty{33.9}{s^{-2}}$：$T = 2\pi/5.83 = \qty{1.08}{s}$。由 $T = \qty{1.10}{s}$：
$\gamma = 4\pi^2mV/(T^2PS^2) = 4\pi^2 \times 0.0165 \times 0.010/(1.21 \times 10^5
\times 4\times10^{-8}) = 1.35$.

\textbf{24.} 摩擦力使运动衰减，若为\omterm{prop:b1:newton-dynamics:friction}{干摩擦}，还会使
平衡位置移动；小球与管壁之间的泄漏使空气在压缩时逸出，
降低回复力：两者都会使周期变长，并使
$\gamma$ 偏低——如上面的 $1.35$。

\textbf{25.} 打气筒与柴油机：$Q = 0$，故 $\Delta U = W$——功转化为
\omterm{def:b1:kinetic-theory:U}{内能}与温度，由拉普拉斯定律（指数为
$\gamma$）可知。小球：快速压缩为\omterm{prop:b1:first-law:transformations}{绝热过程}，且 $\gamma$ 决定
空气弹簧的劲度系数。
\end{solution}
